\section{Conclusion}

As stated in \ref{research_target}, this research aims at the future development of "popular music" and focuses on rhythm. The method of "music with multi-layered temporal structures" conceived in this study can be summarized as follows based on the current results. First, "performance is possible through structural understanding without possessing exceptional technical skills." Second, "this approach can create music that lies within the extension of popular music context while being difficult ensemble pieces." Finally, "production and understanding require the development of notation and production environments." The performers in this study have professional careers, but they do not regularly engage in activities specializing in playing such music. However, in this case, they were able to perform based on their own understanding through lead sheets as meta-structures and their instruction manuals, as well as facilitation that promoted structural understanding. Furthermore, the author succeeded in composing pieces with conventional forms like "Etude \#2" and pieces within the jazz context like "Etude \#1" using the method of "music with multi-layered temporal structures" described in this paper.

The author already had experience learning the Berklee method when creating these pieces, which contributed to the success of the aforementioned compositions. One reason why the Berklee method developed popular music is that "Berklee College of Music did not impose restrictions when spreading its theory." In other words, the theory spread worldwide as musicians who studied at Berklee became educators themselves. Similarly, research on such methods holds great significance when it is made "open access" and "inherited or developed." In this sense, as understood from the interviews in this research, although limited in scope, this study has led other musicians to attempt to update their own musical expression. The author believes this is a greater achievement than the novelty of the research itself, when viewed from the perspective of "thinking about future music from popular music," which is the purpose of this research.\\

Finally, I will discuss future prospects from this research.\\

First, in order to explore the possibilities of rhythm through pieces with such structures, I would like to develop rhythmic cognition and structural understanding of music obtained from this through the development of production environments. This could include harmonic development, co-creation with computers using algorithms and AI, and the possibility of "welcoming such environments as performers or collaborators." Additionally, I would like to continue performance activities with the performers from this study and continue considerations on rhythmic cognition through interviews and collaborations with performers from other genres.

Regarding the development of the concept of "multi-layered temporal structures" proposed in this research, I consider the following. First, expansion into new harmony that could be created from the methods proposed in this research. By considering harmonic development together with multi-layered rhythms using the simple techniques shown in this research, there is potential to create new methods of harmonic development as an alternative to conventional harmonic extensions. As one attempt to demonstrate this possibility, "Etude \#3" was composed by dividing the same chord progression rules into two time signatures. With such a method, triads with simple sonorities can wrap around, and the harmonic sound necessarily becomes more complex/polytonal/atonal as time progresses. In improvisation, performers sometimes "assume" harmony different from the chord progression written in the lead sheet and perform with their own extensions/interpretations, but by conducting such harmonic extension concepts in relation to rhythm, the unique harmonic design that "multi-layered temporal structures" possess could be possible from the composition stage, and if used as an improvisation method, it could be effective as an alternative approach to conventional methods. Also, as a possibility for embodiment, by taking up the relationship with physical expression such as dance, it is necessary to consider not just the viewing experience, as with performance, but the process of how dancers themselves "embody" polyrhythm/cross-rhythm. Through such attempts, there is potential to create relationships not as a master-servant relationship like "dance and its background music," but as a relationship as "performers who create the music listening experience" together, and similarly in relation to other timeline arts such as video works, to create "new listening experiences in bodily sensation." There is still ample room for the development of musical art derived from rhythm.
